\documentclass[../../../include/open-logic-section]{subfiles}

\begin{document}

\olfileid[id]{sth}{ordinals}{ordtype}
\olsection{Ordinal sebagai Tipe Urutan}

Berbekal Penggantian, dan dengan demikian kini bekerja dalam $\ZFminus$,
akhirnya kita dapat membuktikan hasil yang telah kita tuju:

\begin{thm}\ollabel{thmOrdinalRepresentation}
Setiap pengurutan baik isomorfik dengan suatu ordinal yang unik.
\end{thm}

\begin{proof}
Misalkan $\tuple{A, <}$ suatu pengurutan baik. Menurut
\olref[basic]{ordisoidentity}, pengurutan itu isomorfik dengan paling banyak
satu ordinal. Jadi, untuk reductio, andaikan $\tuple{A, <}$ tidak isomorfik
dengan \emph{ordinal mana pun}. Pertama-tama kita akan ``membuat
$\tuple{A, <}$ sekecil mungkin''. Secara terperinci: jika suatu segmen awal
sejati $\tuple{A_a, <_a}$ tidak isomorfik dengan ordinal mana pun, terdapat
$a \in A$ terkecil dengan sifat tersebut; kemudian misalkan $B = A_a$ dan
$\mathord{\lessdot} = \mathord{<_a}$. Jika tidak, misalkan $B = A$ dan
$\mathord{\lessdot} = \mathord{<}$.

Menurut definisi, setiap segmen awal sejati dari $B$ isomorfik dengan suatu
ordinal, yang unik seperti di atas. Jadi, menurut Penggantian, himpunan
berikut ada dan merupakan suatu fungsi:
\[
	f = \Setabs{\tuple{\beta, b}}{b \in B\text{ dan }\ordeq{\beta}{\tuple{B_b, \lessdot_b}}}
\]
Untuk menyelesaikan reductio, kita akan menunjukkan bahwa $f$ merupakan
suatu isomorfisme $\alpha \to B$, untuk suatu ordinal $\alpha$.

Jelas bahwa $\ran{f} = B$. Menurut \olref[iso]{lemordsegments}, $f$
mempertahankan pengurutan, yakni $\gamma \in \beta$ jika dan hanya jika
$f(\gamma) \lessdot f(\beta)$. Untuk menunjukkan bahwa $\dom{f}$ merupakan
suatu ordinal, menurut \olref[basic]{corordtransitiveord} cukup ditunjukkan
bahwa $\dom{f}$ bersifat transitif. Jadi, tetapkan $\beta \in \dom{f}$,
yakni $\ordeq{\beta}{\tuple{B_b, \lessdot_b}}$ untuk suatu $b$. Jika $\gamma
\in \beta$, maka $\gamma \in \dom{f}$ menurut
\olref[iso]{wellordinitialsegment}; dengan menggeneralisasi, $\beta
\subseteq \dom{f}$.
\end{proof}

Hasil ini mengabsahkan definisi berikut, yang ingin kita berikan sejak
\olref[vn]{sec}:

\begin{defn}
Jika $\tuple{A, <}$ merupakan suatu pengurutan baik, maka tipe urutannya,
$\ordtype{A, <}$, adalah ordinal unik $\alpha$ sedemikian sehingga
$\ordeq{\tuple{A, < }}{\alpha}$.
\end{defn}

Selain itu, definisi ini mengabsahkan dua prinsip yang baik:

\begin{cor}\ollabel{ordtypesworklikeyouwant}
Jika $\tuple{A, <}$ dan $\tuple{B, \lessdot}$ merupakan pengurutan baik:
\begin{align*}
	\ordtype{A, <} = \ordtype{B, \lessdot}&\text{ jika dan hanya jika }\ordeq{\tuple{A, <}}{\tuple{B, \lessdot}}\\
	\ordtype{A, <} \in \ordtype{B, \lessdot}&\text{ jika dan hanya jika }\ordeq{\tuple{A, <}}{\tuple{B_b, \lessdot_b}}\text{ untuk suatu }b \in B
\end{align*}
\end{cor}

\begin{proof}
Identitas tersebut berlaku menurut \olref[basic]{ordisoidentity}. Untuk
membuktikan klaim kedua, misalkan $\ordtype{A, <} = \alpha$ dan
$\ordtype{B, \lessdot} = \beta$, dan misalkan $f \colon \beta \to \tuple
{B, \lessdot}$ merupakan isomorfisme kita. Maka:
\begin{align*}
	\alpha \in \beta&\text{ jika dan hanya jika }\funrestrictionto{f}{\alpha} \colon \alpha \to B_{f(\alpha)}\text{ merupakan suatu isomorfisme}\\
	&\text{ jika dan hanya jika }\ordeq{\tuple{A, <}}{\tuple{B_{f(\alpha)}, \lessdot_{f(\alpha)}}}\\
	&\text{ jika dan hanya jika }\ordeq{\tuple{A, <}}{\tuple{B_b, \lessdot_b}}\text{ untuk suatu $b \in B$}
\end{align*}
menurut \olref[iso]{ordisounique},
\olref[iso]{wellordinitialsegment}, dan
\olref[basic]{ordissetofsmallerord}.
\end{proof}

\end{document}
